\chapter{Projeto de automação industrial}
\label{cap:projeto}

\begin{objetivos}
Ao final deste capítulo você deve ser capaz de:
\begin{itemize}
  \item situar a documentação básica dentro do ciclo de um projeto de automação;
  \item ler um P\&ID e extrair dele as informações necessárias ao controle;
  \item montar uma lista de I/O a partir da instrumentação identificada;
  \item reconhecer o que um descritivo operacional precisa conter;
  \item entender como a planta baixa e a arquitetura definem I/O local ou remoto.
\end{itemize}
\end{objetivos}

\section{O projeto antes do programa}

Programar o CLP é a última etapa de um processo que começa muito antes. Um projeto de automação ---
básico, detalhado ou executivo --- se apoia em documentação consistente que define o que o sistema
deve fazer. Sem ela, o programa vira registro de decisões improvisadas, e a manutenção futura fica
dependente de quem escreveu o código.

\begin{figure}[H]
  \centering
  \includegraphics[width=\linewidth]{04-fluxo-projeto.pdf}
  \caption{Fluxo simplificado de um projeto de automação, do requisito à entrega técnica.}
  \label{fig:fluxo-projeto}
\end{figure}

\begin{atencao}
Documentação não é burocracia posterior. É insumo de projeto: a lista de I/O define o porte do
controlador (Capítulo~\ref{cap:porte}), e o descritivo operacional define a lógica
(Capítulo~\ref{cap:boas-praticas}). Documento atrasado é retrabalho garantido.
\end{atencao}

\index{projeto de automação}\index{documentação!de projeto}

\section{Documentos básicos e essenciais}

Cinco documentos formam o núcleo do projeto. Os dois primeiros sustentam os demais:

\begin{table}[H]
  \centering
  \caption{Documentos básicos de um projeto de automação.}
  \label{tab:documentos}
  \begin{tabular}{p{3.8cm} p{6.4cm} p{4.2cm}}
    \toprule
    Documento & O que define & Serve de base para \\
    \midrule
    P\&ID (fluxograma de processo e instrumentação)
      & Equipamentos, tubulação, instrumentação, malhas de controle, intertravamentos e segurança
      & Lista de I/O, descritivo operacional e definição das malhas \\
    Lista de I/O
      & Quantidade e tipo de pontos: DI, DO, AI e AO, por instrumento
      & Dimensionamento de cartões, remotas e porte do controlador \\
    Descritivo operacional
      & Como o processo funciona, o que deve acontecer em cada condição e quais são as lógicas de
        segurança e intertravamento
      & Programação e especificação de modos de operação \\
    Planta baixa de instrumentação
      & Localização geográfica dos instrumentos e o caminho dos cabos
      & Decisão entre I/O local e I/O remoto; arquitetura de rede \\
    Arquitetura de automação
      & Concepção geral do sistema e seus componentes principais
      & Detalhamento e especificação final de compra \\
    \bottomrule
  \end{tabular}
\end{table}

\section{P\&ID: o ponto de partida}

O P\&ID é um diagrama de engenharia de processo que mostra o fluxo físico --- equipamentos,
tubulação e instrumentação --- e é nele que se enxerga o sistema de controle necessário. Uma vez o
P\&ID definido, com a instrumentação exigida pelo funcionamento do processo e suas respectivas
malhas de controle, segurança e intertravamento, inicia-se o \textbf{levantamento de pontos}.

\`A medida que a instrumentação é identificada, cada tag do P\&ID precisa ser analisada quanto ao
sinal que ocupa no controlador. É aqui que nasce a lista de I/O.

\section{Lista de I/O}

A leitura dos sinais de entrada é feita por cartões específicos, e o mesmo vale para as saídas. Os
sinais mais comuns em campo são 4--20\,mA, 1--5\,V CC, 0/24\,V CC, além de entradas para RTD
(PT-100) e termopar em milivolts. Nem todos os tipos cabem no mesmo cartão --- por isso os pontos
são classificados e contados separadamente:

\begin{itemize}
  \item \textbf{DI} (\emph{discrete input}) --- entrada discreta;
  \item \textbf{DO} (\emph{discrete output}) --- saída discreta;
  \item \textbf{AI} (\emph{analog input}) --- entrada analógica;
  \item \textbf{AO} (\emph{analog output}) --- saída analógica.
\end{itemize}

\begin{table}[H]
  \centering
  \caption{Exemplo de lista de I/O.}
  \label{tab:lista-io}
  \begin{tabular}{r l l l c c c c}
    \toprule
    Item & Tag & P\&ID & Serviço & DI & DO & AI & AO \\
    \midrule
    1 & FIT-1003402/1 & DE-001-023 & Vazão de alimentação do R-2371 & 1 &   & 1 &   \\
    2 & XV-1003401/1  & DE-001-023 & Saída de produtos do R-2373      & 2 & 1 &   &   \\
    3 & SX-1003503    & DE-001-023 & Controle de vazão                & 2 & 1 & 1 & 1 \\
    \bottomrule
  \end{tabular}
\end{table}

A soma dos pontos define quantos cartões e quantas remotas serão necessários e, com isso, o porte
do controlador mais adequado ao processo (Capítulo~\ref{cap:porte}).

\begin{pratica}
Preencha a Tabela~\ref{tab:lista-io} com cinco instrumentos de um processo que você conheça. Para
cada um, decida explicitamente:
(a) o sinal é discreto ou analógico?
(b) quantos pontos de cada tipo ele consome?
(c) existe algum ponto que \emph{não} seja de medição nem de comando, mas precise entrar na lista
--- por exemplo, um contato de status de partida de motor ou um alarme de falha do inversor?
Esse último item é o erro mais comum em listas de I/O iniciantes.
\end{pratica}

\section{Descritivo operacional}

O descritivo operacional descreve o funcionamento do processo. É ele que permite entender tanto o
processo quanto a contribuição da automação sobre ele, e é a fonte direta da lógica: condições de
partida, sequência, paradas, alarmes e intertravamentos.

Um descritivo útil responde, para cada equipamento:

\begin{enumerate}
  \item quais são as condições para partir e para parar;
  \item qual é a sequência normal de operação, e o que acontece se um passo não for cumprido;
  \item o que fazer em falha de energia, de comunicação, de instrumento e de equipamento;
  \item quais situações exigem parada de emergência e qual o estado seguro resultante;
  \item o que o operador pode comandar manualmente, e sob quais restrições.
\end{enumerate}

\section{Planta baixa, I/O remoto e arquitetura}

A planta baixa com a localização dos instrumentos define o caminho dos cabos até o sistema de
controle. Dependendo da distribuição e da densidade de instrumentos na área, decide-se entre I/O
local --- todos os sinais indo até o rack do CLP --- e I/O remoto ou distribuído, com as remotas
instaladas no campo.

A instalação com I/O local consome muito cabeamento, borneiras, caixas de passagem, bandejas,
projeto e mão de obra. Os blocos de I/O remotos reduzem esses gastos de forma significativa, porque
os sinais param de percorrer longas distâncias até o painel. A decisão é técnica e financeira, e
resulta na arquitetura de rede do sistema (Capítulo~\ref{cap:redes-scada}).

O documento de \textbf{arquitetura geral de automação} consolida a concepção do sistema: define os
componentes principais que orientarão o detalhamento e a especificação final de compra.

\begin{atencao}
Arquitetura definida sem a planta baixa é chute. Antes de fechar a arquitetura, responda: quantos
instrumentos existem por área? Qual a distância máxima? Quanto custa o metro de cabo, considerando
bandeja, passagem e mão de obra? Qual é o tempo de resposta exigido pelo processo?
\end{atencao}

\section{Fechamento: o que fica para as próximas etapas}

Ao final desta fase o projeto tem: P\&ID definido, lista de I/O fechada com folga, descritivo
operacional escrito, arquitetura definida e a especificação de compra encaminhada. Com isso o
trabalho de programação deixa de ser criativo e passa a ser verificável --- o programa implementa o
que está descrito, e o descritivo permite conferir se implementou.

\begin{exercicios}
\begin{enumerate}[label=\alph*)]
  \item Qual documento define a quantidade de cartões de entrada e saída de um projeto? Justifique.
  \item A uma lista de I/O de um motor foi anotado apenas o comando de partida (DO). Aponte o que
        está faltando e por que isso impediria a lógica de diagnóstico funcionar.
  \item Um P\&ID indica vazão, pressão e nível em uma mesma linha, com uma válvula de controle.
        Descreva a malha e diga quais pontos de I/O ela provoca.
  \item Dê um exemplo de situação em que o I/O remoto é pior escolha que o I/O local no rack.
  \item Escreva, em no máximo dez linhas, o descritivo operacional de uma bomba que deve partir
        manualmente e parar automaticamente por nível alto, com alarme de falha elétrica.
\end{enumerate}
\end{exercicios}
